\scanpage{43}
\margin{42}
Soient $n$ un entier $\ge3$, et considérons un
polygône régulier à \inserted{$m=$} $2n$ côtés, dans un
espace vectoriel \inserted{réel} $V$ de dimension 2.
L'ens des sommets $\widetilde{\unsure{I|J}}$ du polygone forme
\inserted{l'ens des sommets d'}un
polygône combinatoire (on dira aussi qu'il est muni d'une structure de polygone
combinatoire, exprimable de diverses façons équivalentes bien connues), dont
la donnée équivaut, à isom. can. près, à celle du plan vectoriel réel $V$
et du polygône régulier dans $V$. La donnée de l'un ou de l'autre
équivaut en effet à celle d'un torseur sous le groupe diédral
\[
(1)\qquad D_m=\mu_m\,\unsure{\rtimes|\cdot}\,\mathbb Z/2\mathbb Z
\]
\margin{(produit $\frac12$ direct, $\mathbb Z/2\mathbb Z$ opérant sur $\mu_m$ via
$\zeta\mapsto\zeta^{-1}$)}
qui est le groupe d'automorphismes de la structure standard
\[
(2)\qquad V_0=\mathbb C,\quad
\widetilde{\unsure{I|J}}_0=\mu_m\overset{\mathrm{def}}={}_m\mathbb C^*\subset V_0
\]
$\mu_m$ opérant sur $(V_0,\widetilde{\unsure{I|J}}_0)$ par homothéties
\[
z\longmapsto\zeta z\qquad(\zeta\in\mu_m)
\]
et $\mathbb Z/2\mathbb Z$ par conjugaison complexe
$z\longmapsto\overline z$

\margin{il suffirait $n\ge2$, mais le cas $n=2$ \unsure{est un cas}
très dégénéré, \unsure{qui ne présente pas certaines particularités combinatoires},
\unsure{manque bien} d'intérêt, \unsure{car} c'est le cas \unsure{de} d'un trou!}
